\section{Härtung der Demokratie gegen Populismus}
\subsection{Motiviation}

Die Auseinandersetzung mit dem Populismus muss auf mehreren Ebenen erfolgen.
Wirklich besiegen kann man ihn nur, wenn man ihm die inhaltlichen Grundlagen entzieht.

Offenbar ist dies aber nicht so einfach, denn auch in der Demokratie gibt es Elitenbildung und viel schlimmer noch: Abkopplung der Eliten von der Lebensrealität der Nicht-Eliten.
Die Existenz von Eliten an sich ist eher etwas Gutes: Wir alle wollen nicht von Menschen geführt werden, die dümmer sind als wir.
Eigentlich wollen wir auch nicht von Menschen geführt werden, die genauso klug sind wie wir.
Stattdessen wollen wir von den Besten unter uns geführt werden.

Problematisch wird es, wenn immer größere Teile der Bevölkerung den Eindruck gewinnen, dass
\begin{enumerate}[label=\alph*)]
    \item die Eliten in ihrer eigenen Welt leben, die mit der Realität nichts mehr zu tun hat, oder
    \item nur ihre eigenen Interessen verfolgen, die den Interessen des \enquote{Normalbürgers} entgegenlaufen, oder
    \item die Zeiten nicht richtig einschätzen können, oder
    \item einfach auch keine Lösung haben, oder
    \item die Eliten gar nicht wirklich klüger sind als sie, sondern nur privilegiert.
\end{enumerate}

Dem kann man eigentlich nur entgegenwirken, indem man neu Menschen von unten nach oben bringt.
Die Demokratie lässt dies im Gegensatz zur Monarchie ja auch ohne Revolution zu.

Allerdings macht sie es auch nicht unbedingt einfach, von unten nach oben aufzusteigen.
Mit wohlhabenden Eltern in der richtigen Partei hat man schon deutlich bessere Aufstiegschancen als ein Kind aus bildungs- und politikfernen Schichten.

Trotzdem ist es möglich, dass auch Menschen, die das Leben kennen, bis an die Spitze kommen.
Es mag aber dauern, bis es wieder jemand schafft, in eine Führungsposition zu kommen, der wirklich eine gute Führungspersönlichkeit ist, und nicht nur jemand, der gut reden kann oder die notwendige Portion Vitamin B hat.

